\subsection*{Intended Learning Outcomes - Lecture 01}

We want the students to be able to:
\begin{itemize}
\item explain why correlation does not imply causation.
\item explain ~6 correlation inducing scenarios (e.g.\ cause-effect, effect-cause, feedback, 
    confounding, selection bias, deterministic constraints, etc.).
\item write down and explain the definition of causal effect (in real world terms).
\item explain how one can detect causal relationships experimentally.
\item explain randomized controlled trials and argue about advantages and disadvantages.
\item explain the difference between (conditional) observational distributions and interventional distributions, when they are different, when they are the same.
\item analyse real world studies and claims according to the concepts introduced above.
\item work with (elementary) probability theory, write down the main definitions and compute probabilities and expectation values.
\item translate measure theoretic notations, like measure integrals or expectation values, into familiar ones 
    in the discrete and absolute continuous cases.
\end{itemize}
